%% ============================================================================
%% main.tex  --  MOCK Literature Review, used to test lrreview.cls
%%
%% Everything below is invented placeholder content (mock paper, mock data,
%% mock images). Replace the metadata and the sections to write a real review.
%% Build with:  latexmk -xelatex main.tex      (or: xelatex main.tex, twice)
%% ============================================================================
\documentclass{lrreview}

%% ---- Metadata (front page and running headers) ------------------------------
\lrtitle{A Mock Study of Time-Varying Risk Premia in Synthetic Equity Markets}
\lrrunningtitle{A Mock Study of Time-Varying Risk Premia}
\lrauthors{Alex Placeholder and Jordan Sample}
\lryear{2026}
\lrjournal{Journal of Mock Finance, 1(1), 1\textendash 20}
\lrdoi{10.0000/mock.0001}
\lrarea{Asset pricing; mock empirical finance}
\lrseries{Asset Pricing}
\lrpapertype{Mixed (mock empirical and mock theoretical)}
\lrreviewer{Arthur Faugeron}
\lrdatereviewed{3 October 2026}

\begin{document}

%% ---- Front page: metadata table plus one tight abstract paragraph -----------
\lrfrontpage{%
This is placeholder text written to exercise the template. The mock paper asks
whether the premium earned by a synthetic value factor varies with a mock
volatility state, using 240 months of simulated returns for five quintile
portfolios. The mock authors estimate a predictive regression of next-month
excess returns on a standardised state variable, find a positive and
statistically significant slope in the full sample, and show that the result
survives three robustness checks that are described only in outline. The mock
theoretical section derives a two-state premium from a stylised consumption
model. None of the numbers, tables, figures or references in this document are
real: they exist only so that the typography, spacing, tables, equations,
figures, callout blocks and bibliography can be tested and compared with the
HTML and WeasyPrint reports of the library. The abstract is kept to a single
paragraph so that, together with the metadata table, it fits on the first page.}

%% ---- 1. Introduction ---------------------------------------------------------
\section{Introduction}

\subsection{What the paper is about}
\lrtag{Paper} The mock authors argue that the premium on a synthetic value factor
is not constant, and that its variation can be forecast by a single state
variable. This paragraph tests the body face, the justified setting, the inline
tag label and a citation in author-year form, such as Placeholder and Sample
(2026) or (Doe 2019; Roe and Poe 2021).

\lrtag{Critique} A second paragraph, carrying a different tag, checks the vertical
rhythm between paragraphs. The identification strategy is a predictive
regression; the main threat is that the state variable is estimated with error,
which is a mock concern invented for the sake of the test.

\subsection{Executive understanding}
\subsubsection{What is the paper trying to discover?}
Whether a mock state variable forecasts a mock factor premium, once a mock
control is included.

\subsubsection{Why does the question matter?}
It is a stand-in for the second-level headings of a real review. The list
below checks the en dash marker and the spacing between items.
\begin{itemize}
  \item The first item is long enough to wrap onto a second line so that the justified setting and the hanging indent of the marker can be inspected in context.
  \item The second item is short.
  \item The third item contains an inline equation, $r_{t+1} = \alpha + \beta s_t + \varepsilon_{t+1}$, to check that inline math keeps the line spacing even.
\end{itemize}

%% ---- 2. Research problem -----------------------------------------------------
\section{Research Problem and Hypothesis}

\subsection{Research question}
The precise mock question is whether a one-standard-deviation increase in the
state variable $s_t$ raises the expected excess return of the factor by
$\beta$ percentage points per month, after controlling for a mock size effect.

\subsection{Existing literature and gap}
\lrtag{Inference} Two mock strands exist: constant-premium models (Doe 2019) and
time-varying premium models (Roe and Poe 2021). The gap is a mock one: neither
strand uses the synthetic sample described here.

%% ---- 3. Theory ---------------------------------------------------------------
\section{Theoretical Framework and Mechanism}

\subsection{A two-state premium}
\lrtag{Paper} Let the stochastic discount factor be
\begin{equation}
  m_{t+1} = \delta \left( \frac{c_{t+1}}{c_t} \right)^{-\gamma} ,
  \label{eq:sdf}
\end{equation}
and let the expected excess return of an asset satisfy
\begin{equation}
  \mathbb{E}_t\!\left[ R^{e}_{t+1} \right] = - \operatorname{Cov}_t\!\left( m_{t+1},\, R^{e}_{t+1} \right) \big/ \mathbb{E}_t\!\left[ m_{t+1} \right] .
  \label{eq:premium}
\end{equation}
\begin{lrdefs}
  \lrdef{$m_{t+1}$}{The mock stochastic discount factor between dates $t$ and $t+1$.}
  \lrdef{$\gamma$}{Relative risk aversion; larger values make the premium more sensitive to consumption growth.}
  \lrdef{$\delta$}{The subjective discount factor.}
  \lrdef{$R^{e}_{t+1}$}{The excess return of the mock factor portfolio.}
\end{lrdefs}
Equation (\ref{eq:premium}) says that the premium is the negative covariance
of the return with the discount factor, scaled by the mean discount factor;
an asset that pays off in bad times (high marginal utility) earns a lower
premium. A display with a sum and an integral checks the math font:
\begin{equation}
  \hat{\beta} = \frac{\sum_{t=1}^{T-1} \left( s_t - \bar{s} \right) r_{t+1}}{\sum_{t=1}^{T-1} \left( s_t - \bar{s} \right)^{2}} ,
  \qquad
  P(\tau) = \int_{0}^{\infty} \phi(u)\, e^{-\lambda u}\, du .
  \label{eq:ols}
\end{equation}

%% ---- 4. Data and strategy ----------------------------------------------------
\section{Data and Empirical Strategy}

\subsection{Data}
The mock sample has 240 monthly observations of five quintile portfolios.
Table~\ref{tab:data} is a short table, kept whole on one page.

\begin{lrtable}{Mock datasets used in the study}\label{tab:data}
\begin{tabularx}{\linewidth}{@{}L L L r@{}}
\lrtoprule
\lrth{Series} & \lrth{Construction} & \lrth{Sample} & \lrth{Observations}\\
\lrheadrule
Quintile returns & Sorted on a mock value signal & Jan 2006 to Dec 2025 & 240\\
\lrrowrule
State variable & Standardised mock volatility index & Jan 2006 to Dec 2025 & 240\\
\lrrowrule
Controls & Mock size and momentum factors & Jan 2006 to Dec 2025 & 240\\
\lrbottomrule
\end{tabularx}
\lrsource{Source: mock data generated for this template.}
\end{lrtable}

\subsection{Empirical specification}
\lrtag{Paper} The predictive regression is Equation (\ref{eq:ols}) with a control vector.
A figure follows to check the caption, rule and source line.

\lrfigure{assets/mock_timeseries.png}
  {Cumulative return of the mock strategy and the mock benchmark}
  {Source: mock data generated for this template; 240 months; units are percent. The shaded band marks a mock drawdown episode.}

%% ---- 5. Results --------------------------------------------------------------
\section{Results}

\subsection{Main regression}
\lrtag{Evidence} The mock slope is positive and significant. A one-standard-deviation
rise in the state variable corresponds to a mock increase of 0.45 percentage
points in next-month excess return.

\begin{lrtable}{Predictive regressions of next-month excess return on the mock state variable}
\begin{tabularx}{\linewidth}{@{}L r r r r@{}}
\lrtoprule
\lrth{Specification} & \lrth{Slope} & \lrth{Std. error} & \lrth{$t$-statistic} & \lrth{$R^2$ (\%)}\\
\lrheadrule
(1) State variable only & 0.45 & 0.12 & 3.75 & 4.1\\
\lrrowrule
(2) Adding size control & 0.41 & 0.12 & 3.42 & 4.6\\
\lrrowrule
(3) Adding momentum control & 0.38 & 0.13 & 2.92 & 5.0\\
\lrrowrule
(4) Full controls & 0.36 & 0.13 & 2.77 & 5.3\\
\lrbottomrule
\end{tabularx}
\lrsource{Source: mock estimates. Standard errors are mock Newey-West values with six lags.}
\end{lrtable}

\lrfigure{assets/mock_scatter.png}
  {Predictive relation between the mock state variable and next-month excess return}
  {Source: mock data generated for this template; 300 mock observations; the line is an ordinary least squares fit.}

\subsection{Portfolio sorts}
\lrfigure{assets/mock_quintiles.png}
  {Mean monthly return of mock quintile portfolios}
  {Source: mock data generated for this template; error bars are one mock standard error.}

%% ---- 6. Robustness: a long table that may split across pages -----------------
\section{Robustness and Additional Tests}

Table~\ref{tab:robust} is a longer table that is allowed to split across pages
and repeats its header row.

\begin{lrtable}{Robustness checks, concern, test, result}\label{tab:robust}
\lrsans\fontsize{8.4}{10.5}\selectfont
\begin{xltabular}{\linewidth}{@{}L L L L@{}}
\lrtoprule
\lrth{Concern} & \lrth{Test} & \lrth{Result} & \lrth{Resolved?}\\
\lrheadrule
\endhead
\lrbottomrule
\endfoot
Outliers drive the result & Winsorise at 1\% and 99\% & Slope 0.40, still significant & Yes\\
\lrrowrule
Sample period & Split sample in two halves & Slope 0.52 and 0.31 & Partly\\
\lrrowrule
Alternative state variable & Replace with a mock credit spread & Slope 0.29 & Partly\\
\lrrowrule
Look-ahead in the state variable & Lag the variable by one month & Slope 0.33 & Yes\\
\lrrowrule
Multiple testing & Bonferroni over five portfolios & Significant at 5\% for two portfolios & Partly\\
\lrrowrule
Alternative benchmark & Replace the mock benchmark & Similar & Yes\\
\lrrowrule
Placebo & Randomly shuffle the state variable & Slope near zero & Yes\\
\lrrowrule
Heteroskedasticity & Wild bootstrap errors & Wider interval, same sign & Yes\\
\lrrowrule
Small-sample bias & Stambaugh correction & Slope 0.31 & Partly\\
\lrrowrule
Out-of-sample forecasting & Expanding window, 120 months initial & Mock R-squared of 1.2\% & Suggestive\\
\end{xltabular}
\lrsource{Source: mock results written for this template.}
\end{lrtable}

%% ---- 7. Critical evaluation: callout blocks -----------------------------------
\section{Critical Evaluation}

\begin{lrfinding}{Established}
The mock slope is positive in the full sample and in three of four specifications.
This block tests the camel rule at the left, the small-capital label and the
spacing above and below the callout.
\end{lrfinding}

\begin{lrfinding}{Supported but conditional}
The result depends on a mock control. A second sentence ensures the callout wraps
over more than one line, so the left rule must follow the full height.
\end{lrfinding}

\begin{lrfinding}{Suggestive}
The out-of-sample evidence is weak and consistent with the hypothesis only in sign.
\end{lrfinding}

\begin{lrfinding}{Not established}
That the premium is causal; that the finding generalises beyond the mock sample.
\end{lrfinding}

\subsection{Weaknesses}
\subsubsection{Single most important weakness}
\textbf{Problem.} The state variable is a proxy. \textbf{Why it matters.} Error in the
proxy attenuates the slope. \textbf{Potential bias.} Downward. \textbf{How to test.}
Use two independent proxies and compare.

%% ---- 8--10 compressed --------------------------------------------------------
\section{Contradictions and Edge Cases}
\lrtag{Critique} If the discount factor is almost constant (so that $\gamma \to 0$), Equation
(\ref{eq:premium}) gives a premium of zero. At very large $\gamma$ the premium
becomes unreasonably large, which signals a limit of the mock model.

\section{Further Literature}
\begin{lrtable}{Literature connections}
\begin{tabularx}{\linewidth}{@{}L L L L@{}}
\lrtoprule
\lrth{Paper} & \lrth{Shared idea} & \lrth{Relationship} & \lrth{Why it matters}\\
\lrheadrule
Doe (2019) & Constant premium & Preceded; contradicted by mock evidence & Benchmark model\\
\lrrowrule
Roe and Poe (2021) & State-dependent premium & Closest relative & Same predictor family\\
\lrrowrule
Smith (2024) & Alternative data & Extended & Better measurement\\
\lrbottomrule
\end{tabularx}
\lrsource{Source: mock entries.}
\end{lrtable}

\section{Research Gaps}
\subsection{Paper-level gaps}
\textbf{Gap.} Measurement of the state variable. \emph{Why it matters:} it controls
attenuation. \emph{Existing limitation:} a single proxy. \emph{Way forward:} several proxies.

%% ---- 11. Research ideas ------------------------------------------------------
\section{Research Opportunities}

\begin{lridea}{Idea 1. A mock replication with a second proxy}
\lrfield{Research question}{Does the slope survive a second, independently constructed state variable?}
\lrfield{Motivation}{It addresses the main mock weakness.}
\lrfield{Hypothesis}{The slope is similar in sign and smaller in magnitude.}
\lrfield{Data}{Mock monthly returns and two mock proxies.}
\lrfield{Method}{Predictive regressions with Newey-West standard errors.}
\lrfield{Main risk}{The second proxy may be weaker than the first.}
\end{lridea}

\begin{lridea}{Idea 2. A mock out-of-sample test}
\lrfield{Research question}{Does an expanding-window forecast beat the historical mean?}
\lrfield{Motivation}{Predictive significance in sample is not forecasting value.}
\lrfield{Hypothesis}{A small positive out-of-sample R-squared.}
\lrfield{Data}{Same mock data.}
\lrfield{Method}{Campbell and Thompson style comparison.}
\lrfield{Main risk}{Too few observations.}
\end{lridea}

\begin{lrtable}{Ranking of research ideas (mock scores, 1 to 5)}
\begin{tabularx}{\linewidth}{@{}L r r r r@{}}
\lrtoprule
\lrth{Idea} & \lrth{Novelty} & \lrth{Feasibility} & \lrth{Data} & \lrth{Total}\\
\lrheadrule
1. Second proxy & 3 & 5 & 4 & 12\\
\lrrowrule
2. Out-of-sample test & 2 & 4 & 4 & 10\\
\lrbottomrule
\end{tabularx}
\lrsource{Source: mock scores.}
\end{lrtable}

%% ---- 12. Learning framework --------------------------------------------------
\section{Learning Framework}
\subsection{Questions to test understanding}
\subsubsection{Comprehension}
\begin{enumerate}
  \item What does the mock slope mean in units of percentage points per standard deviation?
  \item Why does a proxy attenuate the estimate?
\end{enumerate}

%% ---- 13. Conclusion ----------------------------------------------------------
\section{Conclusion}
\subsection{One-sentence thesis}
A mock state variable forecasts a mock premium, with caveats.
\subsection{Biggest weakness}
The proxy is measured with error.

%% ---- 14. References: hanging indent, unnumbered, no hyperlinks ---------------
\section{References}
\begin{lrbibliography}
\lrref{Doe, J. (2019). A mock constant-premium model. \emph{Journal of Mock Finance}, 1(1), 1\textendash 15.}
\lrref{Placeholder, A., and Sample, J. (2026). A mock study of time-varying risk premia in synthetic equity markets. \emph{Journal of Mock Finance}, 1(1), 1\textendash 20.}
\lrref{Roe, R., and Poe, E. (2021). Mock state-dependent premia. \emph{Mock Review of Financial Studies}, 2(3), 100\textendash 130.}
\lrref{Smith, S. (2024). Mock alternative data for premium estimation. \emph{Working paper}.}
\end{lrbibliography}

\end{document}
